\subsection{All abstract divisor-module morphisms and their full zero fibre}
\label{sec:sdm-all-morphisms}

Retain the original divisors, lattice, space, coefficient sheaf and
modules of SDM1--4. In particular, no coefficient of either divisor
is changed. The divisors are
\[
D=\sum_p n_p[p]+a[\infty],
\qquad
D'=\sum_p n'_p[p]+a'[\infty].
\]
The construction of these modules is that of Alain Connes and
Caterina Consani, \emph{Absolute algebra and Segal's $\Gamma$-sets},
\href{https://arxiv.org/abs/1502.05585v2}{arXiv:1502.05585v2},
Propositions \texttt{propspecz} and \texttt{proppic}, with the full
lattice lift calculated in SDM1--13. The following proof determines
every morphism in this lifted construction, including the maps with
zero amplitude.

Define the subset of the original generic scalar semiring
$G_L(\mathbb Q)$ by
\begin{equation}\tag{SDM14}
\begin{split}
C_L(D,D')={}&\{z_\ell:\ell\in L\}\\
 &{}\cup
\left\{\widehat r:r\in\mathbb Q^\times,\quad
       v_p(r)+n'_p-n_p\ge0\ \text{for every finite }p,\quad
       |r|\le e^{a'-a}\right\}.
\end{split}
\end{equation}
It is a set of permitted scalars; no closure under addition at the
archimedean bound is asserted. For every such scalar, define the
coordinatewise map $f_u$ on every nonempty open and every finite
pointed set by multiplication by $u$. Its exact coordinate formulas
are
\begin{equation}\tag{SDM15}
\begin{aligned}
 f_{\widehat r}(x,\mu)&=(rx,\mu),\\
 f_{z_\ell}(x,\mu)&=(0,\ell\wedge\mu).
\end{aligned}
\end{equation}
In particular the supported-zero scalar $e$ gives
$(x,\mu)\mapsto(0,\mu)$, whereas $\tau$ gives the constant map
to the actual pointed zero. Both occur and are distinct.

\begin{theorem}\label{SDM:all-morphisms}
Multiplication gives the canonical bijection of pointed sets
\begin{equation}\tag{SDM16}
\begin{gathered}
\begin{aligned}
C_L(D,D')&\xrightarrow{\ \sim\ }
\operatorname{Hom}_{\mathcal A_L}
       \bigl(\mathcal M_L(D),\mathcal M_L(D')\bigr),\\
u&\longmapsto f_u.
\end{aligned}\\
C_L(D,D')=\mathcal M_L(D'-D)(X)(1_+).
\end{gathered}
\end{equation}
The equality on the right uses the full labelled-zero fibre.
For a third original divisor $D''$ and
$v\in C_L(D',D'')$, the exact composition law is
\begin{equation}\tag{SDM17}
f_v\circ f_u=f_{vu}.
\end{equation}
Thus multiplication of zero scalars meets their lattice labels;
it does not collapse them to the actual zero.
\end{theorem}

\begin{proof}
First prove that every permitted scalar defines the displayed module
morphism. For $u=\widehat r$ and a local amplitude $x$ satisfying
the source valuation inequalities, one has, for every finite $p$
in the open,
\[
v_p(rx)+n'_p
 =\bigl(v_p(x)+n_p\bigr)
  +\bigl(v_p(r)+n'_p-n_p\bigr)\ge0.
\]
On an open containing infinity, its finite tuple satisfies
\[
\sum_j|rx_j|=|r|\sum_j|x_j|
 \le |r|e^a\le e^{a'}.
\]
These are the original inequalities, with their constants retained.
For $u=z_\ell$, every output amplitude is zero, so every valuation
and positive archimedean bound is satisfied. Formula SDM15 gives its
support exactly.

Multiplication by a fixed scalar commutes with all pointed
pushforwards. For the zero scalar this uses the explicit identity
\[
\ell\wedge\bigvee_{j\in F}\mu_j
 =\bigvee_{j\in F}(\ell\wedge\mu_j)
\]
for every finite fibre $F$, including the empty fibre; its empty
join is $0_L$. For the nonzero scalar it uses distributivity of
rational multiplication over each amplitude sum, with unchanged
support joins. The maps commute with restrictions because those
are inclusions of the same amplitudes and labels. They commute
with the external $\mathcal A_L$-module action by associativity
and commutativity of scalar multiplication and lattice meet.
They therefore define morphisms of the required sheaves of modules.

We now prove that these are every abstract morphism; no preservation
of an amplitude projection or of an individual zero label is assumed.
Let
\[
f:\mathcal M_L(D)\longrightarrow\mathcal M_L(D')
\]
be any morphism of $\mathcal A_L$-module sheaves. The canonical
generic-stalk identifications in SDM4 identify its source, target
and coefficient algebra with
$A_\eta=HG_L(\mathbb Q)$. Write
$u=f_\eta(1_+)(\widehat1)$, where the input is the actual
multiplicative unit $\widehat1=(1,1_L)$ of that generic module.
For every finite pointed set $K$ and every
$x\in A_\eta(K)$, the external action of $x$ on the module
unit at $1_+$ returns the same tuple $x$, using
$K\wedge1_+\simeq K$. Module compatibility therefore gives
\[
f_\eta(K)(x)=x\,u.
\]
This proves the scalar formula at every finite arity, not only
at level one.

For every nonempty open $U$, the map from
$\mathcal M_L(D')(U)(K)$ to its generic stalk is injective:
all restriction maps are coordinatewise inclusions in the common
tuple carrier $G_L(\mathbb Q)^{K^\circ}$, and the stalk is their
filtered union. Naturality with respect to these maps now forces
$f_U(K)$ to be the restriction of multiplication by that same
generic scalar $u$. This proves uniqueness once its permitted
values have been determined.

Every element of $G_L(\mathbb Q)$ is either $z_\ell$ for a unique
$\ell\in L$, or $\widehat r$ for a unique nonzero rational $r$.
The zero-amplitude elements have already been shown to work without
any further inequality. Suppose $u=\widehat r$. At the finite
stalk $p$, apply the map to the literal input
$\widehat{p^{-n_p}}$. It belongs to the source module
$HG_L(p^{-n_p}\mathbb Z_{(p)})$. Its output must belong to
$HG_L(p^{-n'_p}\mathbb Z_{(p)})$, so
\[
v_p(r)-n_p\ge-n'_p.
\]
This is exactly the finite inequality in SDM14.

At infinity, every positive rational $t\le e^a$ is a permitted
level-one amplitude. If $|r|e^a>e^{a'}$, density of the rationals
provides
\[
\frac{e^{a'}}{|r|}<t<e^a.
\]
Its image has absolute amplitude greater than $e^{a'}$, a
contradiction. Hence $|r|\le e^{a'-a}$ is necessary, including
the original endpoint. The sufficiency of precisely these
inequalities was proved at the start. This proves the bijection.
Its basepoint is preserved because $\tau$ acts as the actual
constant-zero morphism.

The equality with global sections in SDM16 is now obtained by
substituting the coefficients $n'_p-n_p$ and $a'-a$ of the
actual difference divisor into SDM2--3. The same full zero lattice
appears on both sides. Finally, composition of the constructed maps
is scalar multiplication by $vu$. If both amplitudes are nonzero,
their finite inequalities add and their archimedean bounds multiply.
If either amplitude is zero, the product is its exact lattice-zero
scalar. In either case $vu\in C_L(D,D'')$, proving SDM17.
\end{proof}

\begin{corollary}\label{SDM:all-isomorphisms}
Every abstract isomorphism of the displayed
$\mathcal A_L$-module sheaves is the scaling in SDM13 and fixes
each labelled zero. More precisely the isomorphisms are exactly
\begin{equation}\tag{SDM18}
\begin{gathered}
f_{\widehat r}\quad(r\in\mathbb Q^\times),\qquad
n'_p-n_p=-v_p(r)\quad\text{for every }p,\qquad
a'-a=\log|r|,\\
D'=D+(r^{-1}),\qquad
f_{\widehat r}(z_\ell)=z_\ell\quad(\ell\in L).
\end{gathered}
\end{equation}
For every original $D$, its complete endomorphism monoid and
automorphism group are
\begin{equation}\tag{SDM19}
\begin{aligned}
\operatorname{End}_{\mathcal A_L}\bigl(\mathcal M_L(D)\bigr)
 &\simeq \{z_\ell:\ell\in L\}\cup
                 \{\widehat1,\widehat{-1}\},\\
\operatorname{Aut}_{\mathcal A_L}\bigl(\mathcal M_L(D)\bigr)
 &=\{f_{\widehat1},f_{\widehat{-1}}\}.
\end{aligned}
\end{equation}
The first isomorphism retains multiplication of the original split
scalars; it does not assert an additive ring structure on this
bounded set.
\end{corollary}

\begin{proof}
An inverse pair induces inverse generic scalar multiplications, by
SDM16--17. A zero-amplitude scalar cannot be invertible, since its
product with any scalar has zero amplitude. A nonzero generic
scalar $\widehat r$ has the unique inverse
$\widehat{r^{-1}}$ in $G_L(\mathbb Q)$. Applying the finite
inequalities of SDM14 to the map and its inverse yields
\[
v_p(r)+n'_p-n_p\ge0,\qquad
-v_p(r)+n_p-n'_p\ge0.
\]
Their left sides sum to zero, so both equal zero. Applying the
two archimedean inequalities gives
\[
|r|\le e^{a'-a},\qquad
|r|^{-1}\le e^{a-a'}.
\]
The second is $|r|\ge e^{a'-a}$; hence equality holds. These are
exactly the conditions in SDM18. Conversely those equalities
put both scalars in the required sets SDM14, so SDM17 proves
that they are inverse sheaf maps. The principal-divisor convention
in SDM12 gives
\[
(r^{-1})=\sum_p(-v_p(r))[p]+\log|r|[\infty],
\]
which proves the claimed identity of divisors with the signs
unchanged. Formula SDM15 proves fixation of every $z_\ell$.

For $D'=D$, the nonzero scalars in SDM14 are rational numbers
with every finite valuation nonnegative and with $|r|\le1$.
Such a rational number is an integer: a reduced denominator
greater than one would have a prime divisor giving a negative
valuation. The only nonzero integers of absolute value at most
one are $1$ and $-1$. All zero scalars remain allowed.
This proves the endomorphism assertion and, by the preceding
inverse calculation, its exact automorphism subgroup.
\end{proof}

\subsection{The complete category of the original divisors}
\label{sec:sdm-divisor-category}

Define a category $\mathsf{Div}_L$ retaining each original divisor
$D$ of SDM1 as an object. Its morphisms, composition and identities
are the specified data
\begin{equation}\tag{SDM20}
\begin{gathered}
\operatorname{Hom}_{\mathsf{Div}_L}(D,D')
   =\mathcal M_L(D'-D)(X)(1_+)=C_L(D,D'),\\
(v,u)\longmapsto v\circ u:=vu,\qquad
\mathrm{id}_D=\widehat1.
\end{gathered}
\end{equation}
The product on the second line is the actual split multiplication:
$\widehat s\,\widehat r=\widehat{sr}$,
$\widehat r\,z_\ell=z_\ell\,\widehat r=z_\ell$ for $r\ne0$,
and $z_\ell z_\mu=z_{\ell\wedge\mu}$.
Every $z_\ell$ is a morphism between every ordered pair of divisors;
no such label has been suppressed.

These data satisfy the category axioms. Closure of composition
is the scalar inequality calculation in the proof of SDM17.
The identity scalar belongs to $C_L(D,D)$ by SDM14, since
its finite valuations are zero and its absolute value is one.
Associativity and the two identity laws follow directly from
associativity and the unit of $G_L(\mathbb Q)$.
The distinguished zero morphism is $\tau$; its products on both
sides are $\tau$. Other $z_\ell$ remain separate morphisms
and have the meet-composition law just displayed.
No additive closure or ring enrichment is asserted: already in
$\operatorname{End}_{\mathsf{Div}_L}(D)$, the amplitude sum
$1+1=2$ fails the archimedean bound in SDM19.

Let $\mathsf{ModDiv}_L$ be the full subcategory of
$\mathcal A_L$-module sheaves with objects the displayed
$\mathcal M_L(D)$ for original divisors $D$. Then
\begin{equation}\tag{SDM21}
\begin{aligned}
\mathcal F:\mathsf{Div}_L&\longrightarrow\mathsf{ModDiv}_L,\\
D&\longmapsto\mathcal M_L(D),\\
u&\longmapsto f_u
\end{aligned}
\end{equation}
is fully faithful and essentially surjective.
Indeed SDM17 proves preservation of composition and SDM15
proves preservation of identities. The bijection SDM16 is
exactly the bijection on every Hom set required for full
faithfulness. Every object of $\mathsf{ModDiv}_L$ is the image
of its displayed original divisor, proving essential
surjectivity. If the full subcategory is enlarged to contain
all module sheaves isomorphic to those displayed objects,
the same argument proves essential surjectivity onto that
enlargement as well: its defining isomorphism is an
isomorphism to an image of $\mathcal F$.

Consequently the original divisor category, with its complete
lattice of zero-amplitude morphisms, is equivalent to the
specified full divisor-module subcategory. The functor retains
the original divisor coefficients and takes each $z_\ell$
to its exact action $(x,\mu)\mapsto(0,\ell\wedge\mu)$.
Its isomorphisms and automorphisms are precisely SDM18--19;
passing to this category does not identify the labelled
zero morphisms or assume an additional purity property.
